%%%
%
% $Autor: Wings $
% $Datum: 2021-05-14 $
% $Pfad: GitLab/CornerBLending $
% $Dateiname: Hints
% $Version: 4620 $
%
% !TeX spellcheck = de_DE/GB
%
%%%


\chapter{Sicherheitsleitfaden und Benutzerhinweise}
\textbf{WICHTIG:} Lesen Sie dieses Handbuch vor der ersten Inbetriebnahme vollständig durch. Bewahren Sie dieses Dokument für die zukünftige Nutzung sorgfältig auf.

\section{Gefahren für den Anwender}
\begin{itemize}
	\item \textbf {Nutzertauglichkeit}
	Der Einfluss von Alkohol, Medikamenten oder Drogen kann zu diversen Gefahren bei der Verwendung des Produkts führen. Daher darf das  Tischförderband nur im nüchternen und gesundheitlich uneingeschränkten Zustand bedient werden.
	\item \textbf {Intaktheit und Funktionsfähigkeit:} 
	Falls das Produkt sichtbare Schäden aufweist, so ist eine weitere Verwendung des Produkts ausdrücklich zu unterlassen. Abhängig vom Produkt und von der Beschädigung besteht nicht nur Verletzungsgefahr (z.B. durch scharfe Kanten), sondern auch Lebensgefahr (z.B. durch einen elektrischen Spannungsübertritt).
	\item \textbf {Aufsicht:}
	Ein Gebrauch durch minderjährige Benutzer in Eigenverantwortung ist nicht gestattet. Unter Aufsicht volljähriger Personen, welche den entsprechenden Sicherheitsanweisungen des Tischförderbands kundig sind, ist eine Benutzung an dieser Stelle zulässig.
	\item \textbf {Einsatzort:}
	Das Tischförderband und das Netzteil sind nur für den Betrieb in trockenen Innenräumen geeignet. Schützen Sie diese vor Feuchtigkeit und berühren Sie diese niemals mit nassen oder fettigen Händen.
	\item \textbf {Kleinteile:}
	Bewahren Sie Verpackungsmaterialien und lose Gegenstände wie das  USB-Kabel außerhalb der Reichweite von Kleinkindern auf (Erstickungs- und Strangulationsgefahr).
\end{itemize}

\bigskip

\section{Mechanische Sicherheit}
\begin{itemize}
	\item \textbf {Stützbeine:}
	Das Tischförderband verfügt über Stützbeine. Prüfen Sie vor jedem Betrieb, ob die Stützbeine vollständig ausgeklappt und fest arretiert sind. Ein instabiler Stand kann zum Umkippen des Tischförderbands oder zum Einklemmen von Gliedmaßen führen.
	\item \textbf {Beschädigungen:}
	Trennen Sie das Netzteil sofort vom Strom, wenn Sie ungewöhnliche Geräusche, Gerüche oder Rauchentwicklung bemerken. Zerlegen Sie das Netzteil nicht. Bei Beschädigungen (z. B. am Stecker oder Gehäuse) muss das Netzteil entsorgt und durch ein Original-Ersatzteil ersetzt werden.
	\item \textbf {Kleidung und Haare:}
	Tragen Sie keine weite Kleidung, Schals oder Schmuck, die vom   Transportband oder den Laufrollen erfasst werden könnten. Langes Haar muss zusammengebunden werden.
\end{itemize}
